\section{Introduction}\label{sec:intro}

A forecasting signal and a trading strategy are different objects. The first is a statement about conditional expectations; the second must also survive the bid--ask spread, the cost and availability of short positions, execution delay, and the selection effects that accompany any search over many specifications. A large literature documents how often the second test fails where the first succeeds. Anomaly returns shrink after realistic trading costs \citep{novymarx2016}, concentrate in the short leg and largely vanish once borrow fees are charged \citep{muravyev2025}, decay after publication \citep{mclean2016}, and are frequently false discoveries once the number of tried specifications is taken into account \citep{harvey2016,harvey2020,chordia2020,hou2020}. Backtests are fragile because the researcher chooses among many variants on the same data \citep{bailey2014notices,bailey2017pbo,novymarx2015}, and because small violations of the decision-time information set inflate measured performance \citep{kaufman2012,zhang2026}.

Order flow is a natural place to look for predictability that is real. Microstructure theory links signed order flow to subsequent price changes \citep{kyle1985}. Empirically, order imbalance predicts returns of individual stocks and of the market at short horizons \citep{chordia2002,chordia2004}; order-flow imbalance across book levels explains contemporaneous price moves within and across stocks \citep{cont2014,cont2023}; and learned summaries of the limit order book carry forecastable information at several horizons \citep{kolm2023}. The same literature hints at why such signals can be hard to monetize: much of the predictability sits at horizons where turnover is large relative to the forecastable move.

This paper is an audit of what survives when both questions, prediction and monetization, are asked of one commercially available dataset under a protocol designed to prevent self-deception. The setting is the Moscow Exchange (MOEX) in 2025--2026. Following the restrictions on non-resident trading introduced in 2022, the market has been dominated by domestic retail investors, who accounted for 70.7\% of equity transactions in 2025 \citep{moex2026retail}. The Bank of Russia key rate ranged between 21\% and 14\% over our sample \citep{cbr2026}, so capital committed to a strategy has a high opportunity cost. Short selling is costly and patchy: only 10 of the 60 less-liquid stocks we checked were available to borrow at the broker. The exchange sells an order-flow dataset, AlgoPack, with five-minute ``super-candle'' statistics on trades, orders and the order book \citep{moexalgo}. As with the exchange analytics product that \citet{akyildirim2021} study for Borsa Istanbul, the practical question is whether a dataset any subscriber can buy contains cross-sectional information about future returns, and whether a subscriber can turn that information into net profit.

A data-provenance problem shapes the answer. AlgoPack's pre-2025 history was regenerated by the exchange after the fact. Every record carries a system timestamp (\texttt{SYSTIME}), and for equities, records written within 30 days of the trade date begin only on 3 January 2025. We use 2020--2024 only for specification and sign checks and issue every verdict on real-time records. The rule is consequential: the effect of the most predictive single feature is about twice as large on the regenerated history as in real time, every quantity we tuned on the regenerated history was more optimistic than its real-time counterpart, and an apparent improvement of a price-based classifier from order-flow inputs turned out to be market timing that disappears once features are demeaned within date.

We find four things. First, a regularized linear composite of 75 daily order-flow features forecasts the cross-section of 10-day abnormal returns on the real-time window with a mean daily rank information coefficient (IC) of 0.0957 (Newey--West $t=5.87$). It exceeds a baseline built on the 10-level book imbalance by 0.037, its permutation placebos sit at zero, its IC is positive in seven of the eight calendar quarters, and it retains 84\% of its IC when the forecast is one day stale. The ranking does not separate from price history. A 12-month momentum rank, which is not among the composite's inputs, reaches IC 0.1072 ($t=3.95$) on the same window and is confirmed on independent 2021--2024 prices (IC 0.038, $t=2.56$); run through the same book it nets at least as much as the composite; the pre-specified contrast between the two books is $+0.019$\% per day with an interval that contains zero; and the 91\% of the composite's ranking power that survives the removal of the momentum span does not convert into a separable net return. What the evidence establishes at the level of ranking is narrower than what it showed a month ago: the composite is not a momentum proxy, but neither is it worth more than one.

Second, monetization is not established, and the input that decides it is no longer assumed. The short-borrow terms of the admitted universe are now measured at the broker: across the 72 names that resolve, the annual rate runs from 20.0\% to 125.0\%, median 25.0\%, quartiles 22.0\% and 36.9\%, and no admitted name is available at the 17\% that the earlier draft used as its base. Carry enters the net return linearly at about 0.0037\pp\ per day per percentage point, so the whole curve follows from one run: the quoted-spread book earns 0.0679\% per day at 17\%, 0.0383\% at the measured median of 25\% ($t=0.89$, interval $[-0.040,+0.122]$), 0.0199\% at 30\%, and $-0.006$\% at the upper quartile. Under the flat 0.15\% per-side tariff of the live account the same book nets 0.0298\% at 17\% and nothing at 25\%. The headline of the previous draft, 0.075\% per day, was computed at the lowest rate in the programme's assumptions; it is an upper bound, not an estimate. The rest of the earlier reading stands: a one-day signal delay cuts the net return by 70\%, a two-day delay turns it negative, and the deflated Sharpe probability \citep{bailey2014dsr} is 0.685 after ten trials and 0.346 after one hundred. The tradable part of the signal decays much faster than its IC.

Third, part of the economics is now measured, part is still assumed, and the mapping remains one choice among many. The fee is measured, availability by date is not, and neither are executions beyond the quoted touch; we separate measured, bounded and assumed inputs explicitly (Section~\ref{sec:uncertainty}). Under the observed variance, detecting a true net return of 0.075\% per day with conventional power would take roughly 1,230 trading days, and demonstrating that the net return is smaller than the policy rate would take about 2,200; the sample cannot settle the economic question in either direction.

Fourth, the composite is one of 42 research directions run under the same protocol. Thirty-three failed their pre-specified criteria, four passed their criteria without earning more than the policy rate, and five remain leads awaiting frozen forward tests. The testing machinery appears calibrated: none of the mechanism-free control variables (lunar phase, planetary positions, geomagnetic activity, calendar-day parity, bitcoin volatility) or permutation placebos passed in the families for which we report them, and in the two largest libraries the number of nominal $p<0.05$ results was at or below the number expected by chance.

The paper's contributions are correspondingly modest and specific. It provides real-time, out-of-sample evidence on order-flow predictability in an emerging, retail-dominated market under a protocol fixed before outcomes were read. It documents and quantifies a provenance hazard in a widely available commercial dataset. It decomposes where the monetization of a real signal is lost and states which economic inputs would have to be measured to settle the question. And it reports a complete ledger of negative results and the pitfalls that changed results, in the spirit of recent reporting standards for forecasting and trading claims \citep{khushi2026}. It does not claim a trading strategy.

Section~\ref{sec:lit} relates the paper to the literature. Section~\ref{sec:data} describes the data, Section~\ref{sec:method} the protocol and methods. Section~\ref{sec:signal} reports the forecasting results and Section~\ref{sec:money} the monetization results. Section~\ref{sec:uncertainty} discusses economic uncertainty and pre-specified extensions, Section~\ref{sec:secondary} summarizes secondary results and leads, Section~\ref{sec:discussion} discusses implications, Section~\ref{sec:limits} lists limitations and Section~\ref{sec:conclusion} concludes. Appendices contain the feature definitions, secondary modules in detail, the ledger of all 42 directions, the catalogue of pitfalls, the calibration of the testing machinery, the frozen forward tests and a reproducibility checklist.
